% Generated draft: requires mathematical and source review.
\documentclass[11pt]{article}
\usepackage[margin=25mm]{geometry}
\usepackage{fontspec}
\setmainfont{Latin Modern Roman}
\usepackage{amsmath,amssymb,mathtools,unicode-math}
\setmathfont{Latin Modern Math}
\usepackage{enumitem,needspace,xurl,hyperref}
\setlength{\parindent}{0pt}
\setlength{\parskip}{5pt}
\setlength{\emergencystretch}{3em}
\newcommand{\N}{\mathbb{N}}
\newcommand{\Z}{\mathbb{Z}}
\newcommand{\Q}{\mathbb{Q}}
\newcommand{\R}{\mathbb{R}}
\newcommand{\C}{\mathbb{C}}
\newcommand{\F}{\mathbb{F}}
\newcommand{\A}{\mathbb{A}}
\newcommand{\PP}{\mathbb{P}}
\newcommand{\E}{\mathbb{E}}
\newcommand{\eps}{\varepsilon}
\newcommand{\dd}{\,\mathrm{d}}
\begin{document}
\section*{UnsolvedMath ID: 30005860}
\textbf{Source problem code:} catalogue entry
\subsection{Definitions and mathematical statement}
The supplied question refers to the IAS minimum-area compactness conjecture, but does not state its hypotheses. The formulation below is conditional on those hypotheses; they cannot be reconstructed from this record alone.

Let \((M^3,g_j)\) be a sequence in the referenced setting. Here \(M^3\) denotes the common three-dimensional manifold appearing in the source. Suppose, along the same subsequence and with the identifications needed for tensor convergence, that
\[
 g_j\longrightarrow g_\infty\quad\text{in }W^{1,p},
 \qquad
 (M^3,g_j)\xrightarrow{\mathrm{SWIF}}\mathcal X_\infty.
\]
The exponent \(p\), the precise Sobolev topology, and the background structures used to compare tensors are not specified. Sobolev regularity concerns tensor coefficients and their first weak derivatives. SWIF denotes Sormani--Wenger intrinsic-flat convergence of the associated integral current spaces; their orientations and integral current structures must be supplied by the referenced setting. Write \(\mathcal X_\infty=(X_\infty,d_\infty,T_\infty)\) when the limit is nonzero. The source describes its underlying metric space as rectifiable and expressly allows the zero integral current space as a limit.

For a specified representative of \(g_\infty\), let
\[
 S=\{x\in M^3:g_\infty(x)\text{ is infinite valued or degenerate}\},
 \qquad U=M^3\setminus S,
\]
as in the source. This description requires further interpretation: a Sobolev tensor is ordinarily defined only almost everywhere, so it does not determine this pointwise set without a representative or another specified convention. Whenever the chosen tensor on \(U\) induces a well-defined finite length metric, denote that metric by \(d_{g_\infty}\). Its length construction is based on
\[
 L_{g_\infty}(\gamma)
 =\int_a^b\sqrt{g_\infty(\gamma(t))(\dot\gamma(t),\dot\gamma(t))}\,\mathrm dt,
 \qquad
 d_{g_\infty}(x,y)=\inf_{\gamma:x\leadsto y}L_{g_\infty}(\gamma).
\]
The admissible curves, the interpretation of the integrand for a rough tensor, and any needed finiteness or separation hypotheses remain to be specified. Write \(\widehat U\) for the metric completion of \((U,d_{g_\infty})\), when this metric is defined.

\textbf{Questions.} For sequences satisfying the referenced IAS hypotheses and the stated Sobolev convergence, must there exist a metric-space isometry
\[
 (X_\infty,d_\infty)\cong\widehat U?
\]
If this identification does not follow from the stated setting, what additional hypotheses guarantee it, and what relation between these spaces can otherwise be established? The source explicitly permits a zero intrinsic-flat limit, so the proposed identification must address that case rather than assume that \(X_\infty\) is nonempty. It also remains necessary to specify whether the comparison uses the underlying set of the integral current space or its metric completion. The question asks for an isometry of metric spaces, without explicitly asking that an isometry preserve an integral current structure.

Finally, suppose that a specified definition of distributional scalar curvature is applicable to \(g_\infty\) and that
\[
 \operatorname{Scal}(g_\infty)\geq0
 \quad\text{in the distributional sense}.
\]
When scalar curvature is realized as a distribution on \(M^3\), this inequality means
\[
 \langle\operatorname{Scal}(g_\infty),\varphi\rangle\geq0
 \quad\text{for every }\varphi\in C_c^\infty(M^3)
 \text{ with }\varphi\geq0,
\]
using the density and pairing prescribed by that definition. What conclusions about \(\mathcal X_\infty\), or about its relation to \(\widehat U\), follow from this inequality? The record does not specify the distributional framework, its regularity assumptions, whether it includes the singular set, or a particular desired conclusion. Consequently this last part is a research question rather than a fully specified implication.
\subsection{Short English statement}
For a three-dimensional sequence in the referenced IAS minimum-area compactness setting, compare its intrinsic-flat limit with the completion of the region where its Sobolev limit metric is finite and nondegenerate. Ask whether these spaces are isometric, which additional assumptions ensure this, and what distributional nonnegative scalar curvature implies for the intrinsic-flat limit. The IAS hypotheses, Sobolev exponent, singular-set convention, and distributional framework are missing from the supplied record.
\subsection{Sources}
\begin{itemize}
\item Local UnsolvedMath registry, supplied complete record: permanent ID \(30005860\), title \emph{Intrinsic-Flat and Sobolev Limit Spaces}, exact source code \texttt{catalogue entry}. The record attributes the question to the Oberwolfach Reports collection, lists proposed year \(2024\), and supplies no individual proposer.
\item \emph{Analysis, Geometry and Topology of Positive Scalar Curvature Metrics} (2024), Oberwolfach Reports, pp.\ 483--566, DOI \texttt{10.4171/owr/2024/9}. Supplied source: \url{https://doi.org/10.4171/owr/2024/9}. The quoted excerpt identifies Open Question 3 and includes the printed page marker \(549\).
\item Supplied literature source, labelled ``original OWR report'': \url{https://doi.org/10.4171/owr/2024/9}.
\item Supplied literature source, labelled ``Sormani--Tian--Wang scalar-curvature limit report'': \url{https://arxiv.org/abs/2404.17121}.
\item Supplied literature source, labelled ``Allen scalar-curvature stability report'': \url{https://arxiv.org/abs/2404.02662}.
\item The supplied literature assessment was checked at \texttt{2026-08-22} and labels the problem open. It states that known examples and partial compactness theorems do not establish the proposed identification, or a general distributional scalar-curvature principle of the requested kind. This is supplied dated curation, not an independently verified literature assessment.
\item The registry labels the statement \texttt{corrected\_verified} and reports that it was reconstructed and checked against the official Oberwolfach report. This is a supplied verification claim, not verification performed for this draft. The excerpt also mentions an upcoming paper without identifying it.
\end{itemize}
\end{document}
